\scanpage{106}
laquelle disparait quand on passe au carré
\[
(\det B)^{\otimes2}\otimes\deleted{\unsure{A/B}}(A/B)^{\otimes2}\simeq1
\]
qui donne
\[
(\det B)^{-\otimes-2}\simeq(A/B)^{\otimes2}
\]
Ainsi, on peut écrire
\[
\delta_{\det\mid B}\in(A/B)^{\otimes2}\simeq D^{\otimes-2}
\]
[et en fait, c'est simplemt donné par la restriction de la forme quadratique $\det$ à $D$ : $\delta_{\det\mid B}=\det_{\mid D}$]

Or on a une application bilin can
\diagram{\begin{tikzcd}[ampersand replacement=\&,row sep=small]
\deleted{\unsure{\mu:}}\ \deleted{\unsure{A}}B\times B \arrow[r] \& A/B\\
(u,v) \arrow[r] \& uv\bmod B
\end{tikzcd}}
et celle ci est \unsure{alterné|alternée}, puisque par H.C, on a pour $u\in B$
\[
u^2\equiv0\qquad\bmod B
\]
(\deleted{\unsure{car}} $u^2=\operatorname{Tr}(u)u+\det u.1$, et $1,u\in B$).

Ainsi, on trouve
\[
\bigwedge^2 B\xrightarrow{\mu} A/B
\]
\unsure{de|Or|On} sorte que
\[
\mu(u\wedge v)=uv\bmod B.
\]
D'ailleurs \deleted{\unsure{$\mu$}} la fonction bil. $uv\bmod B$ est nulle quand l'un ou l'autre des \unsure{argument|arguments} est dans $k$, donc on peut considérer $\mu$ comme
% archive sheet number: 105
